% Related Work --- NEW section (the Avance7 executive summary has none).
% Every citation is anchored to its concrete use in the project; we make
% no claims about papers beyond how they are used here.
\section{Trabajo Relacionado}
\label{sec:related}

\subsection{Modelos de fundación para la observación de la Tierra}
AlphaEarth Foundations~\cite{brown2025alphaearth} provee embeddings
anuales globales listos para análisis, distribuidos en Google Earth
Engine como el activo \texttt{SATELLITE\_EMBEDDING/V1/ANNUAL} (versión de
datos~1.1, 64-dimensional, CC-BY-4.0). Usamos estos embeddings
directamente como características tabulares en lugar de entrenar un
modelo de fundación propio. DINOv3 en su variante
satelital~\cite{oquab2023dinov3} aporta características congeladas
autosupervisadas. AnySat~\cite{astruc2024anysat} es un modelo de
observación de la Tierra multirresolución que evaluamos como backbone de
segmentación. Para la comprensión fina y alineada con el lenguaje,
FarSLIP~\cite{li2025farslip} adapta CLIP a la teledetección y aporta la
rama contrastiva visión-lenguaje usada aquí. Trabajo reciente sobre la
transferibilidad espacial de estos embeddings reporta que el desempeño
se degrada entre regiones y que la adaptación few-shot es
necesaria~\cite{harvesting2026alphaearth}; nuestros experimentos de
transferencia lo corroboran.

\subsection{Segmentación de series temporales de imágenes satelitales}
Los modelos de atención temporal establecieron el estado del arte en el
mapeo de cultivos a nivel de parcela sobre series temporales de
imágenes. U-TAE introdujo un codificador ligero de autoatención temporal
para SITS~\cite{garnot2021utae}, y el benchmark PASTIS / PASTIS-R
formalizó la evaluación panóptica y semántica con folds espacialmente
disjuntos~\cite{garnot2021pastis}. TSViT~\cite{tarasiou2023tsvit}
factoriza la atención a lo largo de los ejes temporal y espacial y es el
backbone más fuerte en nuestro estudio; lo extendemos con una rama
fenológica (TSViT-pheno) inspirada en los transformadores conscientes de
la fenología~\cite{wen2025phenology}. Como baselines puramente
espaciales incluimos U-Net~\cite{ronneberger2015unet},
DeepLabv3+~\cite{chen2018deeplabv3plus} y
SegFormer~\cite{xie2021segformer} (variante B0, RGB de tres bandas).
Señalamos que, a diferencia de planes previos del proyecto, no se
entrenó ningún modelo volumétrico Swin-UNETR;
AnySat~\cite{astruc2024anysat} ocupa el lugar del modelo de fundación en
nuestra comparación.

\subsection{Modelos visión-lenguaje para la agricultura}
Benchmarks multimodales como AgroMind~\cite{ruan2025agromind} evalúan la
comprensión de teledetección agrícola mediante respuesta a preguntas.
AgroMind es solo de evaluación (sin partición de entrenamiento, alrededor
de 28{,}482 pares de QA); el ajuste fino sobre él induciría fuga de
datos, por lo que lo usamos estrictamente para la evaluación de la capa
conversacional. Nos apoyamos en modelos de lenguaje abiertos ---
Gemma~\cite{team2024gemma} y Qwen3~\cite{yang2025qwen3} (ambos
Apache~2.0) --- y en el razonador frontera
Gemini~2.5~Pro~\cite{google2025gemini25} para la operación en la nube.

\subsection{Agentes LLM anclados}
El patrón \emph{Be My Eyes}~\cite{huang2025bemyeyes} desacopla la
percepción del razonamiento: los perceptores emiten observaciones
textuales y un razonador congelado responde, lo que acota la
alucinación porque cada cifra citada se origina en un resultado de
herramienta auditable. Adoptamos este patrón al pie de la letra: nuestros
modelos de segmentación y tabulares son los perceptores, y el LLM es un
razonador congelado que nunca clasifica píxeles. Combinamos además un
prefiltro espacial PostGIS con búsqueda semántica pgvector (Spatial-RAG)
para anclar las respuestas en parcelas vecinas reales.
